\documentclass[11pt]{article}

\usepackage[a4paper,margin=1in]{geometry}
\usepackage{amsmath,amssymb,amsfonts}
\usepackage{bm}
\usepackage{booktabs}
\usepackage{hyperref}
\usepackage{enumitem}
\usepackage{ctex}

\title{LL-Gaussian New Full-Process Decomposition and Loss Design}
\author{LL-Gaussian New Structure Notes}
\date{\today}

\begin{document}

\maketitle

\begin{abstract}
本文档用于整理下一版 LL-Gaussian 的全过程结构设计。它沿用当前项目中
``反射率--光照--瞬时残差--增强光照'' 的低光分解思路，但将主路径重新整理为
\textbf{B0 + offset detail + reflectance decoder} 的反射率表达、
\textbf{SG low-light illumination} 的低光光照表达、
\textbf{SG enhanced illumination} 的增强光照表达，以及
\textbf{single-head residual} 的瞬时误差吸收表达。

本文尤其明确一点：新结构中 residual 保持单头即可，不再把 residual 拆成
noise head 和 artifact head。Residual 的职责是局部、延迟、受 mask 约束地吸收
$R \odot L$ 主分解无法解释的瞬时误差，而不是成为第二套图像生成分支。
\end{abstract}

\section{Design Goal}

新结构的核心目标是让低光重建和增强结果具备更清晰的分工：
\begin{itemize}[leftmargin=1.5em]
    \item $R$ 负责稳定材质颜色、纹理、边缘和局部结构；
    \item $L$ 负责当前低光观测下的视角相关光照变化；
    \item $L^{\text{enh}}$ 负责把同一反射率重光照到增强状态；
    \item $E$ 作为单头 residual，只负责局部困难区域的瞬时补偿；
    \item 最终增强结果默认不引入 residual，避免噪声吸收分支污染正式增强输出。
\end{itemize}

给定一张低光观测图像 $I$，训练期主重建写为
\begin{equation}
    I^{\text{rec}}
    =
    \mathrm{clip}
    \left(
        R \odot L + E_{\text{used}},
        0, 1
    \right),
\end{equation}
其中
\begin{equation}
    E_{\text{used}} = w_{\text{res}}(t) \, M_{\text{res}} \odot E.
\end{equation}

增强结果写为
\begin{equation}
    I^{\text{enh}} = \mathrm{clip}\left(R \odot L^{\text{enh}}, 0, 1\right).
\end{equation}

因此，主路线可以概括为
\begin{equation}
    \boxed{
    I \approx R \odot L + w_{\text{res}} M_{\text{res}} \odot E,
    \qquad
    I^{\text{enh}} \approx R \odot L^{\text{enh}}
    }
\end{equation}

\section{Full Process Overview}

新结构的完整流程分为七个阶段：
\begin{enumerate}[leftmargin=1.8em]
    \item 数据读取与相机、点云、depth prior 初始化；
    \item Gaussian anchor 与 offset 表达初始化；
    \item 反射率 $R$、低光光照 $L$、增强光照 $L^{\text{enh}}$、单头 residual $E$ 的网络参数初始化；
    \item 每次迭代根据当前相机执行 differentiable splatting，得到多张渲染图；
    \item 根据训练阶段组合 reconstruction、SG、reflectance、enhancement、residual、geometry 等损失；
    \item 按 warmup、normal、residual ramp、enhancement guidance 等时间调度逐步开放分支；
    \item 训练结束后渲染 train/test/interp，并分别输出低光重建、增强结果、诊断图和指标。
\end{enumerate}

\section{Scene Representation}

\subsection{Anchor and Offset}

场景仍由 Gaussian anchors 和局部 offsets 表达。第 $i$ 个 anchor 具有：
\begin{itemize}[leftmargin=1.5em]
    \item anchor position $\mathbf{x}_i$；
    \item anchor feature $\mathbf{f}_i$；
    \item $K$ 个 offset $\mathbf{o}_{i,k}$；
    \item opacity、scaling、rotation 等 Gaussian 参数；
    \item 反射率底色 $B_i^0$；
    \item offset 级反射率 detail $\Delta^R_{i,k}$；
    \item 增强光照 SG 参数；
    \item 可选 residual 相关 feature、offset、scaling 参数。
\end{itemize}

一次渲染中，只有当前视角可见且通过 opacity/filter 选择的 anchor-offset 会参与 rasterization。
因此本文中的 $R$、$L$、$L^{\text{enh}}$、$E$ 都先在 offset 级生成，再通过 splatting 投影到图像空间。

\subsection{Initialization}

初始化建议保持当前代码中的三层策略：
\begin{itemize}[leftmargin=1.5em]
    \item anchor 和 offset 来自输入点云；
    \item $B^0$ 使用第一帧或代表帧做 Retinex 风格初始化；
    \item offset detail 初始化为零，让初始反射率主要由 $B^0$ 控制；
    \item reflectance decoder 初始化为接近零修正；
    \item SG illumination 和 enhanced SG illumination 初始化为稳定、低能量、不过尖锐的状态；
    \item residual 初始化为接近零输出，且训练早期不参与主重建。
\end{itemize}

\section{Reflectance Branch}

\subsection{Three-Level Reflectance}

新结构保留当前较有效的三层反射率设计。

\paragraph{(1) Anchor-level base reflectance.}
每个 anchor 学习一个三通道对数反射率：
\begin{equation}
    B_i^0 \in \mathbb{R}^3,
    \qquad
    R_i^{\text{base}} = \exp(B_i^0).
\end{equation}

$B^0$ 负责材质底色和低频颜色稳定性，不期望吸收明显高光、噪声或局部曝光异常。

\paragraph{(2) Offset-level log-detail.}
每个 offset 学习一个有界 detail：
\begin{equation}
    \Delta^R_{i,k} \in \mathbb{R}^3.
\end{equation}

在 log-reflectance 空间中加入 detail：
\begin{equation}
    R_{i,k}^{\text{detail}}
    =
    \exp\left(
        B_i^0 + s_R \tanh(\Delta^R_{i,k})
    \right),
\end{equation}
其中 $s_R$ 是 detail scale。该分支用于恢复局部纹理、边缘和小尺度材质变化。

\paragraph{(3) View-independent decoder refinement.}
在 $B^0 + \Delta^R$ 之后，使用一个不接 view direction 的轻量 decoder 做局部微调：
\begin{equation}
    D^R_{i,k}
    =
    f_{\theta_R}\left(\mathbf{f}_i, \mathbf{o}_{i,k}\right).
\end{equation}

最终 offset 级反射率写为
\begin{equation}
    R_{i,k}
    =
    \mathrm{clip}
    \left(
        R_{i,k}^{\text{detail}}
        \odot
        \left(1 + \gamma_R \tanh(D^R_{i,k})\right),
        \epsilon_R,
        1
    \right).
\end{equation}

这里 decoder 只作为局部 refine，不应替代 $B^0$ 和 offset detail 成为自由颜色生成器。

\subsection{Reflectance Responsibility}

反射率分支的职责边界如下：
\begin{itemize}[leftmargin=1.5em]
    \item 允许学习物体固有颜色、纹理、边缘、可重复出现的细节；
    \item 不应学习视角相关照明；
    \item 不应吸收随机噪声；
    \item 不应吸收只在少量视角出现的彩色高光或曝光异常；
    \item 不应为了匹配增强伪 GT 而被长期拉平。
\end{itemize}

\section{Low-Light Illumination Branch}

\subsection{SG Illumination}

低光光照使用 spherical Gaussian 表达。对每个 offset，SG branch 输出 $M$ 个 lobes：
\begin{equation}
    \left\{
        \left(
            \bm{\mu}_{i,k,m},
            \lambda_{i,k,m},
            \mathbf{a}_{i,k,m}
        \right)
    \right\}_{m=1}^{M}.
\end{equation}

其中：
\begin{itemize}[leftmargin=1.5em]
    \item $\bm{\mu}_{i,k,m}$ 为单位方向；
    \item $\lambda_{i,k,m}$ 为 sharpness；
    \item $\mathbf{a}_{i,k,m}$ 为 RGB amplitude；
    \item 当前视角方向记为 $\mathbf{v}_{i}$。
\end{itemize}

低光光照为
\begin{equation}
    L_{i,k}
    =
    \frac{1}{M}
    \sum_{m=1}^{M}
    \mathbf{a}_{i,k,m}
    \exp
    \left(
        \lambda_{i,k,m}
        \left(
            \bm{\mu}_{i,k,m}^{\top}\mathbf{v}_{i} - 1
        \right)
    \right).
\end{equation}

\subsection{Illumination Responsibility}

$L$ 只解释当前低光图像中的光照和曝光状态：
\begin{itemize}[leftmargin=1.5em]
    \item 允许随视角和局部几何变化；
    \item 允许表达低光照明不均、暗角、方向性照明；
    \item 不应承担纹理锐化的主要职责；
    \item 不应把材质颜色完整吸收进去；
    \item 不应为了重建误差而退化成高频彩色贴图。
\end{itemize}

\section{Enhanced Illumination Branch}

\subsection{Enhanced SG}

增强分支仍坚持 ``同一反射率，不同光照'' 的重光照原则。增强光照记为
\begin{equation}
    L^{\text{enh}}_{i,k}.
\end{equation}

建议使用 anchor/offset 级增强 SG 参数与轻量上下文特征共同生成增强光照：
\begin{equation}
    L^{\text{enh}}_{i,k}
    =
    h_{\theta_E}
    \left(
        \mathrm{stopgrad}(\mathbf{f}_i),
        \mathrm{stopgrad}(\phi^L_{i,k}),
        \Theta^{\text{enh-sg}}_{i,k}
    \right),
\end{equation}
其中 $\phi^L_{i,k}$ 是低光 illumination branch 的中间特征，$\Theta^{\text{enh-sg}}$ 是增强 SG 参数。

增强图像为
\begin{equation}
    I^{\text{enh}} = R \odot L^{\text{enh}}.
\end{equation}

\subsection{Enhanced Branch Responsibility}

增强分支的职责是：
\begin{itemize}[leftmargin=1.5em]
    \item 提升全局与局部亮度；
    \item 保持颜色自然，避免绿色偏色和通道漂移；
    \item 尽量不改变 $R$ 中的材质身份；
    \item 不依赖 residual 输出正式增强结果；
    \item 只通过有限的 pseudo target guidance 获得增强方向，而不是完全复制伪 GT。
\end{itemize}

\section{Single-Head Residual Branch}

\subsection{Why Single Head}

新结构中 residual 明确保持单头：
\begin{equation}
    E_{i,k} = r_{\theta_E}\left(\mathbf{f}_i, \mathbf{o}_{i,k}, \psi_{i,k}\right),
\end{equation}
其中 $\psi_{i,k}$ 包含 view direction、distance、appearance embedding 等当前 residual 输入特征。

不再引入 noise head 和 artifact head 的原因是：
\begin{itemize}[leftmargin=1.5em]
    \item 当前主问题不是 residual 类型不够细，而是 residual 和 $R/L$ 的职责边界要更稳；
    \item 单头 residual 更容易调度、可视化和正则；
    \item 单头可以避免两个 residual head 之间互相补偿、互相逃逸；
    \item 单头与当前最终渲染 ``默认不带 residual'' 的策略更一致。
\end{itemize}

\subsection{Single-Head Residual Equation}

单头 residual 的训练期使用形式为
\begin{equation}
    E_{\text{used}}
    =
    w_{\text{res}}(t)
    M_{\text{res}}
    \odot
    E.
\end{equation}

其中：
\begin{itemize}[leftmargin=1.5em]
    \item $E$ 是 residual net 的原始输出；
    \item $M_{\text{res}}$ 是 residual hard mask；
    \item $w_{\text{res}}(t)$ 是 residual ramp 权重；
    \item residual 未启用时 $E_{\text{used}}=0$。
\end{itemize}

推荐 residual 输出使用有界激活，例如 $\tanh$：
\begin{equation}
    E = \eta_E \tanh(\hat{E}),
\end{equation}
其中 $\eta_E$ 控制最大补偿幅度。若沿用当前代码的 sigmoid 单头兼容路径，也需要通过正则和 ramp 防止其成为全局正向补光器。

\subsection{Residual Mask}

单头 residual mask 建议由两类困难区域合成：
\begin{equation}
    M_{\text{res}} = M_{\text{err}} \cup M_{\text{highlight}}.
\end{equation}

其中：
\begin{align}
    M_{\text{err}}
    &=
    \mathrm{TopK}
    \left(
        \left\|I - \mathrm{stopgrad}(R \odot L)\right\|_1,
        p_{\text{err}}
    \right), \\
    M_{\text{highlight}}
    &=
    \mathrm{TopK}
    \left(
        S_{\text{chroma-highlight}}(I, R),
        p_{\text{hl}}
    \right).
\end{align}

这样 residual 主要在高重建误差区域、彩色高光或局部异常区域工作，不参与整图低频拟合。

\subsection{Residual Scheduling}

Residual 不应从训练一开始参与主重建。建议使用：
\begin{equation}
    w_{\text{res}}(t)
    =
    \begin{cases}
        0, & t < t_{\text{start}}, \\
        \min\left(1, \frac{t - t_{\text{start}} + 1}{T_{\text{ramp}}}\right), & t \ge t_{\text{start}}.
    \end{cases}
\end{equation}

其中 $t_{\text{start}}$ 控制 residual 开启时间，$T_{\text{ramp}}$ 控制线性爬坡长度。该调度保证 $R \odot L$ 先承担主体结构，然后 residual 再处理局部困难误差。

\section{Rendering Equations}

\subsection{Training Render}

训练期对当前视角得到：
\begin{align}
    I^{\text{base}} &= R \odot L, \\
    I^{\text{rec}}  &= \mathrm{clip}\left(I^{\text{base}} + w_{\text{res}} M_{\text{res}} \odot E, 0, 1\right), \\
    I^{\text{enh}}  &= \mathrm{clip}\left(R \odot L^{\text{enh}}, 0, 1\right).
\end{align}

其中 $I^{\text{rec}}$ 用于低光监督，$I^{\text{enh}}$ 用于增强监督和指标。

\subsection{Final Render}

最终常规低光渲染建议输出：
\begin{equation}
    I^{\text{render}} = R \odot L.
\end{equation}

最终增强渲染建议输出：
\begin{equation}
    I^{\text{render-enh}} = R \odot L^{\text{enh}}.
\end{equation}

Residual 默认不参与最终输出。若需要诊断 residual，可额外导出：
\begin{equation}
    I^{\text{diagnostic-res}} = E,
    \qquad
    I^{\text{diagnostic-used}} = w_{\text{res}}M_{\text{res}}\odot E.
\end{equation}

\section{Loss Design}

\subsection{Main Reconstruction Loss}

主重建监督使用 $I^{\text{rec}}$ 对齐低光 GT $I$：
\begin{equation}
    \mathcal{L}_{\text{rec}}
    =
    (1-\lambda_{\text{ssim}})
    \mathcal{L}_{1+}
    \left(
        I^{\text{rec}}, I
    \right)
    +
    \lambda_{\text{ssim}}
    \left(
        1 - \mathrm{SSIM}(I^{\text{rec}}, I)
    \right).
\end{equation}

其中 $\mathcal{L}_{1+}$ 是亮度加权 L1，用于强化暗区误差：
\begin{equation}
    \mathcal{L}_{1+}
    =
    \left|
        \frac{
            I^{\text{rec}} - I
        }{
            \alpha \, \mathrm{stopgrad}(I^{\text{rec}}) + \phi
        }
    \right|.
\end{equation}

\subsection{Illumination and Geometry Loss}

低光 illumination 仍使用像素约束和平滑约束：
\begin{equation}
    \mathcal{L}_{L}
    =
    \lambda_{\text{illu}}
    \mathcal{L}_{\text{illu}}(I, L)
    +
    \lambda_{\text{smooth}}
    \mathcal{L}_{\text{smooth}}(L, I).
\end{equation}

几何部分保留 scaling 和 depth prior：
\begin{equation}
    \mathcal{L}_{\text{geo}}
    =
    \lambda_{\text{scale}}
    \mathcal{L}_{\text{scale}}
    +
    \lambda_{\text{depth}}
    \mathcal{L}_{\text{depth}}.
\end{equation}

\subsection{SG Regularization}

SG illumination 需要限制能量和 sharpness：
\begin{equation}
    \mathcal{L}_{\text{sg}}
    =
    \lambda_{\text{sg-energy}}
    \mathcal{L}_{\text{sg-energy}}
    +
    \lambda_{\text{sg-sharp}}
    \mathcal{L}_{\text{sg-sharp}}.
\end{equation}

该项避免 SG 变成极尖、极不稳定的方向性贴图，同时维持低光照明的可解释性。

\subsection{Reflectance Loss}

反射率损失分为稳定、结构、异常抑制和参数幅度四类：
\begin{align}
    \mathcal{L}_{R}
    ={}&
    \lambda_{\text{cons}}
    \mathcal{L}_{\text{cons}}(R)
    +
    \lambda_{\text{smooth-R}}
    \mathcal{L}_{\text{smooth-R}}(R, L)
    \nonumber \\
    &+
    \lambda_{\text{edge}}
    \mathcal{L}_{\text{edge}}(R, I)
    +
    \lambda_{\text{edge-up}}
    \mathcal{L}_{\text{edge-up}}(R, I)
    \nonumber \\
    &+
    \lambda_{\text{contrast}}
    \mathcal{L}_{\text{contrast}}(R, I)
    +
    \lambda_{\text{hf}}
    \mathcal{L}_{\text{hf}}(R, I)
    \nonumber \\
    &+
    \lambda_{\text{highlight}}
    \mathcal{L}_{\text{highlight}}(R)
    +
    \lambda_{\text{detail}}
    \mathcal{L}_{\text{detail}}
    +
    \lambda_{\text{decoder}}
    \mathcal{L}_{\text{decoder}}
    +
    \lambda_{\text{B0}}
    \mathcal{L}_{\text{B0-3D}}.
\end{align}

各项职责为：
\begin{itemize}[leftmargin=1.5em]
    \item $\mathcal{L}_{\text{cons}}$：抑制反射率中的随机噪声；
    \item $\mathcal{L}_{\text{smooth-R}}$：在 illumination 权重下做轻量平滑；
    \item $\mathcal{L}_{\text{edge}}$：让 $R$ 的结构边缘对齐输入图像结构；
    \item $\mathcal{L}_{\text{edge-up}}$：只惩罚 $R$ 边缘弱于目标结构的位置；
    \item $\mathcal{L}_{\text{contrast}}$：鼓励局部灰度对比；
    \item $\mathcal{L}_{\text{hf}}$：鼓励 Laplacian 高频结构；
    \item $\mathcal{L}_{\text{highlight}}$：抑制反射率吸收彩色亮斑；
    \item $\mathcal{L}_{\text{detail}}$：约束 offset detail 幅度；
    \item $\mathcal{L}_{\text{decoder}}$：约束 decoder refinement 幅度；
    \item $\mathcal{L}_{\text{B0-3D}}$：约束 3D 邻近 anchors 的底色平滑。
\end{itemize}

\subsection{Single-Head Residual Loss}

单头 residual 的损失建议写为：
\begin{align}
    \mathcal{L}_{\text{res}}
    ={}&
    \lambda_{\text{res-sparse}}
    \mathbb{E}
    \left[
        \left|M_{\text{res}}\odot E\right|
    \right]
    +
    \lambda_{\text{res-scale}}
    \mathcal{L}_{\text{scale-res}}
    \nonumber \\
    &+
    \lambda_{\text{res-zero}}
    \left\|
        \mathrm{mean}(M_{\text{res}}\odot E)
    \right\|_1
    +
    \lambda_{\text{res-chroma}}
    \mathcal{L}_{\text{res-chroma}}(E, R).
\end{align}

其中：
\begin{itemize}[leftmargin=1.5em]
    \item sparse 项限制 residual 总幅度；
    \item scale 项限制 residual Gaussian 的体积；
    \item zero-mean 项防止 residual 变成全局补光；
    \item chroma 项只作为轻量引导，把局部彩色异常从 $R$ 中迁移出来；
    \item 不再区分 noise sparse、artifact sparse、noise high-frequency 等双头专属项。
\end{itemize}

如果训练中发现 residual 过度接管重建，应优先调整：
\begin{itemize}[leftmargin=1.5em]
    \item 推迟 $t_{\text{start}}$；
    \item 拉长 $T_{\text{ramp}}$；
    \item 提高 $\lambda_{\text{res-sparse}}$；
    \item 降低 mask 覆盖率；
    \item 降低 residual 输出幅度 $\eta_E$。
\end{itemize}

\subsection{Enhancement Loss}

增强监督由亮度提升、平滑、颜色统计和 refined target guidance 组成。

\paragraph{(1) Brightness degree.}
\begin{equation}
    \mathcal{L}_{\text{degree}}
    =
    \lambda_{\text{local}}
    \mathcal{L}_{\text{degree-local}}
    +
    \lambda_{\text{global}}
    \mathcal{L}_{\text{degree-global}}.
\end{equation}

\paragraph{(2) Enhanced smoothness.}
\begin{equation}
    \mathcal{L}_{\text{enh-smooth}}
    =
    \mathcal{L}_{\text{smooth}}
    \left(
        \frac{L^{\text{enh}}}{\rho},
        I
    \right),
\end{equation}
其中 $\rho$ 是增强倍率或局部亮度参考。

\paragraph{(3) Color statistic alignment.}
对增强结果 $I^{\text{enh}}$ 和 refined target $\hat{I}$ 施加弱颜色统计约束：
\begin{align}
    \mathcal{L}_{\text{color-mean}}
    &=
    \left\|
        \mu(I^{\text{enh}}) - \mu(\hat{I})
    \right\|_1, \\
    \mathcal{L}_{\text{color-std}}
    &=
    \left\|
        \sigma(I^{\text{enh}}) - \sigma(\hat{I})
    \right\|_1.
\end{align}

\paragraph{(4) Green bias penalty.}
\begin{equation}
    \mathcal{L}_{\text{green}}
    =
    \mathbb{E}
    \left[
        \max
        \left(
            0,
            I^{\text{enh}}_G
            -
            \frac{I^{\text{enh}}_R + I^{\text{enh}}_B}{2}
            -
            \tau_G
        \right)
    \right].
\end{equation}

\paragraph{(5) Refined target guidance.}
在指定迭代之后加入 refined target guidance：
\begin{align}
    \mathcal{L}_{\text{diff-illu}}
    &=
    \left\|
        L^{\text{enh}} \odot \mathrm{stopgrad}(R)
        -
        \hat{I}
    \right\|_1, \\
    \mathcal{L}_{\text{diff-refl}}
    &=
    \left\|
        \mathrm{stopgrad}(L^{\text{enh}}) \odot R
        -
        \hat{I}
    \right\|_1.
\end{align}

总 guidance 为
\begin{equation}
    \mathcal{L}_{\text{diff}}
    =
    w_{\text{diff}}(t)
    \left(
        \mathcal{L}_{\text{diff-illu}}
        +
        \lambda_{\text{diff-R}}
        \mathcal{L}_{\text{diff-refl}}
    \right).
\end{equation}

其中 $\lambda_{\text{diff-R}}$ 应保持较小，避免伪 GT 长期把 $R$ 拉平；$w_{\text{diff}}(t)$ 可随训练进度衰减。

\section{Training Schedule}

\subsection{Warmup Stage}

Warmup 阶段目标是先稳定 $R \odot L$：
\begin{itemize}[leftmargin=1.5em]
    \item 使用主重建、illumination、geometry、轻量 reflectance 结构项；
    \item 不启用 residual；
    \item 不启用强 refined target guidance；
    \item 不让增强分支主导整体优化。
\end{itemize}

Warmup 总损失可写为
\begin{equation}
    \mathcal{L}_{\text{warmup}}
    =
    \mathcal{L}_{\text{rec}}
    +
    \mathcal{L}_{L}
    +
    \mathcal{L}_{\text{geo}}
    +
    \mathcal{L}_{R}.
\end{equation}

\subsection{Normal Stage Before Residual}

进入 normal 训练后，但 residual 尚未开启时：
\begin{equation}
    I^{\text{rec}} = R \odot L.
\end{equation}

该阶段继续强化 $R$ 的结构清晰度和 $L$ 的光照解释，并逐步引入 SG regularization 与 enhancement loss。

\subsection{Residual Ramp Stage}

当 $t \ge t_{\text{start}}$ 后，启用单头 residual：
\begin{equation}
    I^{\text{rec}} = R \odot L + w_{\text{res}}(t)M_{\text{res}}\odot E.
\end{equation}

该阶段重点监控：
\begin{itemize}[leftmargin=1.5em]
    \item residual mix weight；
    \item residual mask coverage；
    \item residual abs mean；
    \item residual chroma mean；
    \item reflectance highfreq、edge uplift 和 decoder mean。
\end{itemize}

\subsection{Enhancement Guidance Stage}

当 refined target 可用且迭代达到指定阈值后，加入 $\mathcal{L}_{\text{diff}}$。该 guidance 应更偏向训练 $L^{\text{enh}}$，而不是直接强拉 $R$：
\begin{equation}
    \lambda_{\text{diff-R}} \ll 1.
\end{equation}

\subsection{Total Objective}

完整目标函数为
\begin{align}
    \mathcal{L}_{\text{total}}
    ={}&
    \mathcal{L}_{\text{rec}}
    +
    \mathcal{L}_{L}
    +
    \mathcal{L}_{\text{geo}}
    +
    \mathcal{L}_{\text{sg}}
    +
    \mathcal{L}_{R}
    \nonumber \\
    &+
    \mathcal{L}_{\text{res}}
    +
    \mathcal{L}_{\text{degree}}
    +
    \lambda_{\text{enh-smooth}}
    \mathcal{L}_{\text{enh-smooth}}
    +
    \mathcal{L}_{\text{diff}}
    \nonumber \\
    &+
    \lambda_{\text{mean}}
    \mathcal{L}_{\text{color-mean}}
    +
    \lambda_{\text{std}}
    \mathcal{L}_{\text{color-std}}
    +
    \lambda_{\text{green}}
    \mathcal{L}_{\text{green}}.
\end{align}

\section{Logging and Diagnostics}

建议保留以下日志项：
\begin{itemize}[leftmargin=1.5em]
    \item reconstruction：loss、Ll1、SSIM、PSNR；
    \item illumination：illumination mean、SG energy、SG lambda mean；
    \item reflectance：edge、edge uplift、contrast、highfreq、highlight、detail、decoder；
    \item residual：residual enabled、mix weight、mask coverage、abs mean、raw image、used image；
    \item enhancement：degree local/global、color mean/std、green bias、diff illumination、diff reflectance；
    \item geometry：depth similarity、scaling regularization、visible count。
\end{itemize}

Residual 诊断图建议分三类：
\begin{itemize}[leftmargin=1.5em]
    \item raw residual：$E$ 的真实输出；
    \item used residual：$w_{\text{res}}M_{\text{res}}\odot E$；
    \item normalized abs residual：仅用于观察结构，不代表真实贡献量级。
\end{itemize}

\section{Rendering and Evaluation}

\subsection{Low-Light Test}

低光测试使用
\begin{equation}
    I^{\text{render}} = R \odot L.
\end{equation}

指标文件建议为：
\begin{itemize}[leftmargin=1.5em]
    \item \texttt{metrics\_lowlight.json}：对测试视角原始低光 GT；
    \item 保存 PSNR、SSIM、LPIPS 的 summary 和 per-view 结果。
\end{itemize}

\subsection{Enhanced Test}

增强测试使用
\begin{equation}
    I^{\text{render-enh}} = R \odot L^{\text{enh}}.
\end{equation}

指标文件建议为：
\begin{itemize}[leftmargin=1.5em]
    \item \texttt{metrics\_enhanced\_gt.json}：对数据集 \texttt{gt/images}；
    \item 不把 residual 加入增强指标；
    \item 可额外输出 residual 诊断图，但不参与正式增强评价。
\end{itemize}

\subsection{Interp Video}

插值视频路径没有真实 GT，对应策略为：
\begin{itemize}[leftmargin=1.5em]
    \item 只导出低光重建、增强结果和可选 depth；
    \item 不计算 GT 指标；
    \item 视频封装使用稳定的 OpenCV \texttt{VideoWriter}。
\end{itemize}

\section{Implementation Checklist}

若把本文档作为下一步代码修改参考，建议按以下顺序落地：
\begin{enumerate}[leftmargin=1.8em]
    \item 确认命令行默认 \texttt{use\_dual\_transient=False}，训练脚本不再打开双头 residual；
    \item 保留 \texttt{residual\_net} 单头，逐步弱化或移除 \texttt{noise\_net/artifact\_net} 主路径依赖；
    \item 将 residual mask 统一为单个 \texttt{residual\_mask}；
    \item 将 residual loss 统一为 sparse、zero-mean、scale、chroma 四类；
    \item wandb 中保留 \texttt{residual\_image\_raw} 和 \texttt{residual\_image}，删除双头专属日志；
    \item final render 默认不加 residual，仅保留 \texttt{include\_residual\_render} 作为诊断开关；
    \item 检查 checkpoint/PLY 兼容逻辑，旧 artifact residual 可映射到新 single-head residual；
    \item 重新跑 8k 诊断训练，重点观察 $R$ 清晰度、residual coverage 和增强颜色稳定性；
    \item 再跑 30k 正式训练，比较 \texttt{metrics\_lowlight.json} 与 \texttt{metrics\_enhanced\_gt.json}。
\end{enumerate}

\section{Practical Interpretation}

新结构可以理解为：
\begin{itemize}[leftmargin=1.5em]
    \item $R$ 是 ``可复用的场景材质层''；
    \item $L$ 是 ``当前低光拍摄条件''；
    \item $L^{\text{enh}}$ 是 ``目标增强照明条件''；
    \item $E$ 是 ``训练期局部误差缓冲器''；
    \item final enhanced output 只依赖 $R \odot L^{\text{enh}}$。
\end{itemize}

因此，single-head residual 不是退化设计，而是为了让分解边界更清楚：
它只给 $R \odot L$ 留一个局部缓冲空间，不额外引入一套难以解释的双头瞬时分解。

\section{Open Directions}

后续可继续探索：
\begin{itemize}[leftmargin=1.5em]
    \item 更稳的 residual mask 覆盖率调度；
    \item 更直接的增强结构监督，减少对 refined pseudo target 的长期依赖；
    \item 更物理一致的 enhanced SG 初始化；
    \item 更干净的 checkpoint 兼容路径，把旧双头参数自动迁移到单头 residual；
    \item 让 $R$ 的高频结构提升与 residual 的局部补偿进一步解耦。
\end{itemize}

\end{document}
